\begin{enumerate}[label=\chinese*、,leftmargin=2em,labelsep=0pt,itemsep=8pt,topsep=0pt]
  \item 硕士生的开题报告内容应包括文献综述、选题意义、研究内容、研究方案，论文工作量的估计、工作条件，预期达到的水平，存在的问题及拟采取的解决措施。
  \item 博士生的开题报告内容应包括文献综述、选题背景及其意义、研究内容、研究特色、工作难点、预期成果及其可能的创新点。
  \item 正式开题之前，研究生应在广泛阅读中、外文资料的基础上，深入了解拟选课题的国内外研究动态，把握所选课题的目的、意义和预期结果，明确课题工作的设想、方法和研究路径。
  \item 博士生在规定的时间内，写出开题报告初稿，经指导教师审阅同意后，进行开题报告的送审。开题报告评审未通过者，修后重新送审，若第二次送审仍通不过者，则按有关规定终止学籍。
  \item 不能按期申请开题评议者，应及时向学院提出延期申请。
  \item 本表若不够填写时，可另加附页。
\end{enumerate}
